% 02_explanation(72).tex
\begin{explanationbox}

	\textbf{\textcolor{CExplanation}{一句话直觉}}

	只需知道椭圆的半长轴 $a$、半短轴 $b$ 和直线的倾斜角 $\theta$，直接代入公式即可得到焦点弦长，无需联立方程组。

	\medskip

	\textbf{\textcolor{CExplanation}{核心拆解}}

	\begin{description}[style=nextline, leftmargin=2.8em, labelsep=0.5em, itemsep=0.45em]
		\item[\textbf{要点一（直线参数化）：}]
		      过焦点的直线用倾斜角 $\theta$ 表示参数方程，代入椭圆方程。

		\item[\textbf{要点二（韦达定理）：}]
		      利用韦达定理求出两个交点对应的参数差，从而得到焦点弦长公式
		      \[
			      L=\frac{2ab^2}{a^2-c^2\cos^2\theta}.
		      \]
	\end{description}

	\medskip

	\textbf{\textcolor{CExplanation}{几何本质（方向决定弦长）}}

	焦点弦长 $L$ 依赖于 $a$、$b$ 和 $\cos^2\theta$。

	当 $\theta=90^\circ$ 时，$\cos^2\theta=0$，分母最大，弦长最短：
	\[
		L=\frac{2b^2}{a},
	\]
	这就是椭圆的通径长。

	当 $\theta=0^\circ$ 时，$\cos^2\theta=1$，分母最小，弦长最长：
	\[
		L=2a.
	\]

	\medskip

	\textbf{\textcolor{CExplanation}{代数意义（三角分式）}}

	公式
	\[
		L=\frac{2ab^2}{a^2-c^2\cos^2\theta}
	\]
	可以看成关于 $\cos^2\theta$ 的分式函数。因为
	\[
		0\leq \cos^2\theta \leq 1,
	\]
	所以分母在
	\[
		b^2 \leq a^2-c^2\cos^2\theta \leq a^2
	\]
	之间变化，从而弦长在
	\[
		\frac{2b^2}{a} \leq L \leq 2a
	\]
	之间变化。

	\medskip

	\textbf{\textcolor{CExplanation}{考点价值（弦长与最值）}}

	\begin{itemize}[leftmargin=*, itemsep=0.5em, topsep=0.4em]
		\item \textbf{考法一（直接代入）：} 已知椭圆方程和焦点弦倾斜角，直接代入公式计算弦长。
		\item \textbf{考法二（最值问题）：} 由 $\cos^2\theta$ 的取值范围 $[0,1]$ 求弦长最值，常结合离心率。
		\item \textbf{考法三（反求参数）：} 已知焦点弦长，反求倾斜角、离心率或椭圆参数。
	\end{itemize}

	\medskip

	\textbf{\textcolor{CExplanation}{顿悟点}}

	这类题的关键不是盲目联立，而是看到“过焦点的弦”后，想到用焦点作为起点建立参数方程。这样两个交点对应的参数差，正好就是弦长。

	\medskip

	\textbf{\textcolor{CExplanation}{使用场景}}

	\begin{itemize}[leftmargin=*, itemsep=0.5em, topsep=0.4em]
		\item \textbf{场景一（直接求弦长）：} 题目给出倾斜角或斜率，求焦点弦长。
		\item \textbf{场景二（求最值范围）：} 题目要求焦点弦长的最大值、最小值或取值范围。
		\item \textbf{场景三（反求椭圆参数）：} 已知焦点弦长和角度信息，反推出 $a,b,c$ 或离心率。
	\end{itemize}

\end{explanationbox}